\section{Introduction}

It appears evident that more affluent households live in larger, better houses and more desirable neighborhoods. However, while some characteristics of a neighborhood are based on geography (e.g., waterfronts, elevation) and infrastructure (e.g., commuting time to the workplace), most features that make a neighborhood more or less attractive are a result of the people living in that neighborhood: which shops and restaurants supply the local population? How good is the quality of the local school? What is the crime rate in the neighborhood? Sociologists have therefore described neighborhood formation as an archetype of a \textit{self-organizing process} \citep{schellingDynamicModelsSegregation1971, galsterNatureNeighbourhood2001, fossettEthnicPreferencesSocial2006}. Neighborhoods form from the people that inhabit an area but also shape their daily surroundings and whom they interact with to an extent that shapes their opportunities for life \citep{sampsonAssessingNeighborhoodEffects2002, sharkeyWhereWhenWhy2014, chettyImpactsNeighborhoodsIntergenerational2018a} and subsequent generations \citep{sharkeyIntergenerationalTransmissionContext2008}. As neighborhoods change, they become more or less attractive to both potential new residents and existing residents. The resulting residential mobility can reproduce features of the neighborhood or massively change its local culture and social structure when change reaches a tipping point.

Existing theoretical models of residential segregation and neighborhood formation have explained these phenomena based on local amenities and the preferences of people who they (or do not) want to be their neighbors. Such preferences, combined with income and wealth inequality, lead households with different socio-economic characteristics to sort themselves into distinct neighborhoods, while poorer households are priced out of the most desirable neighborhoods. These models often (implicitly) assume that the type and quality of housing follow such developments. Nevertheless, there is reason to doubt that these models tell the whole picture. Empirical research indicates that households are primarily concerned with finding a place that meets their needs and is affordable. Characteristics of housing units are more important to them than characteristics of the neighborhood \citep{chauCriticalReviewLiterature2003, bayerEquilibriumModelSorting2004, mummoloWhyPartisansNot2017, delucaNotJustLateral2020}. Moreover, changes in the housing stock are a defining feature of neighborhood change \citep{redfernWhatMakesGentrification2003} and precarious housing plays an important role in why the same neighborhoods high shares of in-poverty residents \citep{pendallWhyHighPovertyNeighborhoods2016}. Are the characteristics of the housing units and the agency of those who can change these characteristics (i.e., owners and landlords) not important in explaining residential segregation, neighborhood change, or patterns of inequality? I am not the first to observe this omission. For example, \citet[1081]{rosenthalChangePersistenceEconomic2015} note that "there are important dynamic links between maintenance and neighborhood change although that relationship has mostly escaped attention in the literature" and \citet[1093]{desmondPoorPayMore2019} constitute that landlords are "[c]onspicuously absent from most accounts of neighborhood dynamics or inequality."

Any description of a market is guaranteed to be incomplete if it only discusses one type of market actor, in this case, demand. While the supply side of the housing market, like landlords and real estate agents, has recently gained attention in sociological research in its own right (e.g., \cite{korver-glennCompoundingInequalitiesHow2018, desmondPoorPayMore2019, rosenRacialDiscriminationHousing2021}), this paper aims to combine supply and demand in a theoretical model to explain key workings and patterns in the fields of residential segregation, neighborhoods, residential mobility and housing inequality. The model additionally aims to synthesize and unite these related but previously often unconnected fields \citep{hwangThingsChangeThings2016}. My theoretical model makes three substantial assumptions: households prefer high-quality housing units and prefer to live next to high-status neighbors. Landlords invest in improving the quality of their housing if average rents in the neighborhood increase. If landlords do not invest, the quality of housing decays over time. The model explains the emergence of distinct neighborhoods through a self-reinforcing dynamic: when a neighborhood increases in status, it also becomes more desirable, leading to higher rents and increased investment by landlords. Because the attractiveness of the neighborhood rises with status and housing quality, rents increase further after the investment in housing quality, and poor, typically low-status residents are replaced with more affluent households. This cycle continues until a spatial equilibrium is reached, which aligns well with the reviewed descriptions of gentrification. In equilibrium, the model reproduces several stylized facts from urban research, generating stable and segregated neighborhoods by income, social status, and housing quality. Low-income residents exhibit significantly higher residential mobility, as they are frequently forced to relocate due to affordability issues. They live in lower-quality housing units and neighborhoods, and pay a larger share of their income for housing. As households with low social status reduce the attractiveness of a neighborhood by themselves, they typically live in worse housing conditions than higher-status households with the same income, because landlords invest based on local developments in rents.

The following section summarizes the empirical literature on residential segregation, neighborhoods, and housing inequality to identify the stylized facts that the model should be able to explain. The third section discusses existing theories and their shortcomings, which motivate the development of this new theory. Section four describes the micro-level assumptions of my theoretical model, motivating them with empirical literature on household preferences and real estate investment. Section five describes the dynamics and equilibrium results that these assumptions create in realistic conditions. Before concluding in section seven, I discuss the omissions and limitations of my model in section six.